\section{Related Work and Positioning}\label{sec:related}

\paragraph{CPT optimization and reinforcement learning.}
Prospect theory and \CPT replace expected-utility evaluation by a reference-dependent value transformation and nonlinear probability weighting \citep{kahneman1979prospect,tversky1992advances}.  In sequential decision making, \citet{prashanth2016cumulative} showed that a \CPT criterion must be treated as a functional of the return distribution rather than a transformed immediate reward, and \citet{jie2018stochastic} developed stochastic-optimization machinery for this class of objectives.  The closest reference is \citet{lepel2026policygradients}, whose v5/TMLR formulation derives a first-order \CPT policy-gradient theorem, an order-statistics Monte Carlo estimator, and convergence and sample-complexity guarantees for approximate first-order stationarity.  In parallel, recent optimization work has developed finite-scenario and portfolio-specific algorithms for static \CPT objectives, including ADMM-based decision optimization \citep{cui2025decision}, convex-analytic portfolio procedures \citep{luxenberg2024portfolio}, and a recent behavioral portfolio formulation that combines \CPT with a flexible investment horizon \citep{wu2026behavioral}.  Our work keeps the first-order \CPT policy-gradient perspective but changes two structural elements: the reference becomes an endogenous state generated by realized wealth, and the learning objective therefore becomes a moving, path-dependent sequence.  This requires both exact finite-sample debiasing of the frozen-target gradient and a tracking analysis for the resulting moving objective.

\paragraph{Adaptive references, path dependence, and investor behavior.}
Reference adaptation has long been documented in security-trading experiments, including asymmetric updating after equal-sized gains and losses \citep{arkes2008reference,arkes2010cross}.  Mathematical-finance models then showed that adaptive references can amplify realization behavior \citep{he2019realization}, while path-dependent reference points make current decisions depend on the route by which wealth arrived at its present level \citep{qin2022realization}.  More recent portfolio analysis embeds a dynamic reference directly in multi-period loss-averse allocation and finds materially different policies and disposition-effect predictions \citep{gao2024multiperiod}.  Complementary evidence shows that the disposition effect depends on portfolio-wide gain/loss frames \citep{an2024portfolio}, investor horizon and experience \citep{fan2024horizon}, and the information environment of Chinese retail investors \citep{chen2025xueqiu}.  These studies establish that the relevant behavioral state is neither fixed nor purely asset-local.  They do not, however, provide an online \CPT policy-gradient procedure whose parameter update and first-order tracking guarantees explicitly accommodate reference drift.  That is the point at which the behavioral-reference literature and the optimization problem studied here meet.

\paragraph{Nonstationary and nonconvex online optimization.}
For nonconvex objectives, global comparator regret is generally too strong to characterize first-order learning.  Local-regret criteria instead measure whether recent gradients or gradient mappings are small \citep{hazan2017efficient}.  \citet{aydore2019dynamic} introduced dynamic local regret for nonconvex forecasting with time-smoothed gradients, and \citet{hallak2021regret} extended local-regret ideas to stochastic constrained problems through a proximal-gradient mapping.  In nonstationary RL, variation budgets similarly quantify when a changing environment remains learnable \citep{feng2023nonstationary}.  Recent projected-gradient analysis sharpens first-order stationarity guarantees for constrained nonconvex and stochastic smooth optimization \citep{lan2026projected}.  Our setting combines these ingredients but introduces a distinct source of variation: part of the objective drift is generated by the learner's own realized wealth and reference dynamics.  We therefore work with a normalized projected residual, derive DLR from that residual, and isolate objective variation and simulator error in the online bound.

\paragraph{Debiasing, inference, and information limits.}
The statistical difficulty in a \CPT gradient is generated by a nonlinear functional of an unknown survival distribution.  The order-statistics approach of \citet{lepel2026policygradients} provides a practical first-order estimator and consistency guarantees.  Our estimator targets a different finite-sample property: exact unbiasedness for the regularized frozen-target gradient under a complete-trajectory oracle.  Its independent randomized correction follows the general principle of unbiased randomized telescoping \citep{rhee2015unbiased}, while the leave-one-out expansion uses uniform empirical-distribution control of the type underlying the Dvoretzky--Kiefer--Wolfowitz inequality \citep{massart1990tight}.  The central-limit and minimax layers connect the estimator to classical asymptotic and statistical-decision theory \citep{vandervaart1998asymptotic,tsybakov2009introduction}.  The resulting contribution is therefore not merely a numerical bias correction: it specifies the mean, variance, limiting distribution, expected trajectory cost, and fixed-dimensional information rate of the gradient oracle used by the online algorithm.

\paragraph{Portfolio learning and A-share evaluation.}
The portfolio policy combines factor information with a simplex-valued action.  Dirichlet factor policies provide a natural precedent for reinforced portfolio allocation because they generate valid long-only weights while linking concentration parameters to firm characteristics \citep{andre2020dirichlet}.  For the empirical market layer, recent work emphasizes both the time variation of financial-RL environments \citep{liu2024dynamicdatasets} and the importance of market-specific factor structure in Chinese A-shares \citep{hanauer2024factor}.  These considerations motivate separating the controlled semi-synthetic environment used for theorem-aligned E2--E4 from a chronological real-market test.  Our E5 implements that separation on a fixed 300-stock HS300 universe with strictly point-in-time factors, common transaction costs, and the same stochastic Dirichlet execution rule used by the learner.  The resulting return, drawdown, turnover, and cost outcomes are therefore an external-validity layer for the portfolio mechanism rather than a surrogate test of the DLR theorem.

\paragraph{Human-facing financial agents and heterogeneous risk preferences.}
Recent AI research increasingly treats behavioral preferences as part of the system objective rather than as evaluation noise.  Prospect-theoretic optimization has been used to construct human-aware alignment losses for language models \citep{ethayarajh2024kto}, and recent work directly evaluates whether LLMs represent persona-specific economic risk preferences \citep{liu2025riskpreferences}.  In multi-agent learning, \citet{lalmohammed2025modeling} embed \CPT transformations into heterogeneous agents and show that risk profiles can materially change strategic behavior.  Financial-advice studies add an interaction layer: personalized generative-AI advisors can influence investment choices, but preference elicitation, advice quality, satisfaction, and trust need not move together \citep{takayanagi2025financialadvisors}, while investor resistance to robo-advisors is shaped by both functional and psychological barriers \citep{verma2025roboadvisors}.  Our behavioral design separates two empirical questions that this literature suggests should not be conflated.  E6 holds investor decisions, controls, chronological split, and predictive model fixed while replacing only the reference representation, so its evidence concerns held-out behavioral alignment.  E7 then calibrates heterogeneous investor types from training-period behavior and replays the E5 recommendations with adoption decided before episode returns are observed, separating recommendation uptake from satisfaction after adoption.  This connects dynamic-reference portfolio learning to human-facing interaction without treating simulated satisfaction as observed behavior.